\documentclass[a4paper,10pt, twoside]{report} % Format A4, Schriftgröße 10, zweiseitige Ausrichtung
\usepackage[utf8]{inputenc}
\usepackage[ngerman]{babel} % Sprache: Deutsch nach neuer Rechtschreibung
\usepackage[T1]{fontenc}
\usepackage{lmodern} % andere Schriftart
\usepackage{geometry}
\usepackage{float}
\usepackage{amsmath,amsthm,amsfonts,amssymb} % Mathematik
\usepackage{color,graphicx,overpic}
\usepackage{listings} % Codefragmente
\usepackage{fancyhdr} % Kopf- und Fußzeilen
\usepackage[dvipsnames]{xcolor}
\usepackage{booktabs}
\usepackage[hidelinks,pdfencoding=auto]{hyperref} % Links & Verlinkung
\usepackage{longtable}
\usepackage{cite}
\usepackage[section]{placeins} % Bilder auf der richtigen Seite darstellen
\usepackage{enumitem}
\usepackage{subfiles}
\usepackage{parskip}
\setlength{\headheight}{28.7pt}
\setcounter{secnumdepth}{3}

\pdfinfo{
    /Title (SWP2021 AntiDoS - Entwurfsdokumentation)
    /Creator (TeX)
    /Producer (pdfTeX 3.14)
    /Author (Fabienne Göpfert, Leon Leisten, Felix Hußlein, Robert Jeutter, Jakob Lerch, Tim Häußler, Johannes Lang, Tobias Scholz)
    /Subject (Abwehr von Denial-of-Service-Angriffen durch effiziente User-Space Paketverarbeitung)
}

%% Kopf- und Fußzeilen
\pagestyle{fancy}
\fancyhf{}
\fancyhead[EL]{\leftmark}
\fancyhead[OL]{\includegraphics[height=25pt]{img/aegis_logo_with_name.pdf}}
\fancyhead[OR]{\rightmark}
\fancyfoot[EL]{\thepage}
\fancyfoot[ER]{SWP2021 AntiDoS}
\fancyfoot[OL]{Zweites Reviewdokument}
\fancyfoot[OR]{\thepage}
\renewcommand{\headrulewidth}{0.2pt} %fügt eine Linie unter der Kopfzeile ein

\renewcommand{\arraystretch}{1.3} %sorgt für größeren Abstand zwischen Zeilen in Tabellen

\begin{document}

\begin{titlepage} %Beginn der Titelseite
    \begin{center} %zentriert
        Technische Universität Ilmenau\\
        Fakultät Informatik und Automatisierung \\
        Fachgebiet Telematik/Rechnernetze\\
                %Technische Universität Ilmenau\\
        \vspace{2em}
        \includegraphics[width=0.35\linewidth]{img/tu_ilmenau_logo.png}\\
        %Technische Universität Ilmenau\\
        \vspace{2em}
        Zweite Review zum Thema\\
        \vspace{1em}
        \Large\textsc{Abwehr von Denial-of-Service-Angriffen \\durch effiziente User-Space Paketverarbeitung}\\
        \vspace{0.5em}
        \Large\textsc{Codename: AEGIS}\\
        \vspace{0.5cm}
        \includegraphics[width=0.35\linewidth]{img/aegis_logo.pdf}\\
        \normalsize
        \vspace{2em}
        Autoren:\\
        \begin{table}[h]
            \centering
            \begin{tabular}{l r}
                Fabienne Göpfert & Tim Häußler    \\
                Felix Hußlein   & Robert Jeutter  \\
                Johannes Lang    &  Leon Leisten \\
                 Jakob Lerch   & Tobias Scholz  \\
            \end{tabular}
        \end{table}
        \vspace{1em}
        \normalsize {Betreuer:} Martin Backhaus\\
        \vspace{1em}
        Entwurf 23. Juni 2021
    \end{center}
\end{titlepage}
\pagebreak
\tableofcontents {\protect\thispagestyle{fancy}

\subfile{chapters/1-einleitung}
\subfile{chapters/2-grobentwurf}
\subfile{chapters/3-feinentwurf}
\subfile{chapters/4-bugreview}
\subfile{chapters/5-kimai}
\subfile{chapters/6-abkuerzungsverzeichnis}
\subfile{chapters/7-glossar}


\bibliography{references}
\bibliographystyle{ieeetr}

\listoffigures {\protect\thispagestyle{fancy} % Abbildungsverzeichnis

\end{document}
